% Physical unit 026; accepted Session 03 after fresh 400-dpi authority replay.
\GordanUnitStart{026}{printed page 21}
\GordanTopFolio{21}
\begingroup\small
or, referring it to the whole genus:
\[
 \omega+4n\mathcal A(l_0)-2h\pi,
\]
in which formula $n$ successively assumes the values of all integers, and $h$ represents that integer which is contained in
\[
 {\psi_2+4n\mathcal A(l_0)\over 2\pi}.
\]
Therefore, since
\[
 {\mathcal A(l_0)\over 2\pi}
\]
is irrational (for the contrary case, which occurs very rarely, must be omitted), $\psi$ successively takes all values lying between $0$ and $2\pi$; this was to be demonstrated.

\begin{center}\S\ 4.\end{center}

But now, before we undertake to ask which is the shortest of all lines between two points $A(l\psi)$ and $B(l_1\psi_1)$ on the spheroid, it is useful to consider some matters bearing on that question.

In that question let us put $l$ greater than $l_1$. Since it is now clearly apparent from the preceding theorem that, among the lines issuing from the point $A$, all contain the point $B$, it remains to ask more accurately which exists as shorter than the others. In this case it is indeed evident that the line which is shorter than the others can be defined once its genus is known.

\begin{center}ARGUMENT.\end{center}

For if we consider the course of some geodetic line from the point $A$, and of many other geodetic lines running near it and inclined less toward the meridian of the point $A$, we find that the geodetic line is cut by some line all the sooner, the smaller the inclination toward it is. From this it follows that all geodetic lines of the same genus cut it beyond the terminus. But since every geodetic line, as was shown above, is shortest up to the terminus, all the more, when it cuts another geodetic line, must it long since have ceased to be the shortest line.

\begin{center}\S\ 5.\end{center}

Before, however, we come nearer to determining the genus, let us propose these two theorems.

I. ``The geodetic line $l_0$ I, if, having started from the point $A$ and not yet having reached the point $B$, it cuts the meridian of the point $A$ (for example in $C$), is not the shortest of all lines between $A$ and $B$.''

\begin{center}ARGUMENT.\end{center}

The geodetic line II, whose negative azimuth is equal to the azimuth of the geodetic line I at $A$, by symmetry both touches the circle $l_0$ just as line I does and contains
\endgroup
